协变旋量用于散射外线，而实际束流由普通概率归一化的正能波包制备。两种归一化必须在进入跃迁核之前换算，因为概率投影、自旋求和与相空间使用的内积不同。

\begin{definition}[自由哈密顿量与概率归一化旋量]
令 $\beta=\gamma^0$、$\alpha_i=\gamma^0\gamma^i$，则
\begin{align}
 \{\alpha_i,\alpha_j\}&=2\delta_{ij}I_4,\qquad
 \{\alpha_i,\beta\}=0,\qquad \beta^2=I_4,
 \label{eq:feqo-dirac-anticommutators}\\
 H_0(\vb{p})&=c\boldsymbol\alpha\cdot\vb{p}+\beta m_{\mathrm e}c^2,
 \qquad H_0^2=E_{\vb{p}}^2I_4.
 \label{eq:feqo-dirac-free}
\end{align}
这里 $\alpha_i$ 与 $\beta$ 都厄米，$E_{\vb{p}}>0$。协变式和哈密顿量式之间的词典是
\begin{equation}
 \gamma^0=\beta,\qquad \gamma^i=\beta\alpha_i.
 \label{eq:feqo-gamma-dictionary}
\end{equation}
将\eqnrefs{eq:explicit-u} 的旋量缩放为 $U_s(\vb{p})=\sqrt{c/(2E_{\vb{p}})}\,u_s(p)$，便有
\begin{equation}
 U_s(\vb{p})=N_{\vb{p}}
 \begin{pmatrix}\xi_s\\ C_{\vb{p}}\xi_s\end{pmatrix},
 \qquad
 N_{\vb{p}}=\sqrt{\frac{E_{\vb{p}}+m_{\mathrm e}c^2}{2E_{\vb{p}}}},
 \qquad
 C_{\vb{p}}=\frac{c\boldsymbol\sigma_{\rm P}\cdot\vb{p}}
 {E_{\vb{p}}+m_{\mathrm e}c^2}.
 \label{eq:feqo-positive-spinor}
\end{equation}
大写 $U_s$ 在本段表示单位概率范数旋量；小写 $u_s(p)$ 始终保留\eqnrefs{eq:explicit-u} 的协变归一化。
\end{definition}

\begin{proposition}[同一动量上的正交能量分解]
在固定 $\vb{p}$ 上，定义
\begin{equation}
 \Pi_\pm(\vb{p})=\frac12\left(I_4\pm\frac{H_0(\vb{p})}{E_{\vb{p}}}\right).
 \label{eq:feqo-hamiltonian-energy-projectors}
\end{equation}
则 $\Pi_\pm^2=\Pi_\pm=\Pi_\pm^\dagger$、$\Pi_+\Pi_-=0$、$\Pi_++\Pi_-=I_4$，且 $\Tr\Pi_\pm=2$、$H_0\Pi_\pm=\pm E_{\vb{p}}\Pi_\pm$。同一三动量上的负能本征旋量可取
\begin{equation}
 W_s(\vb{p})=N_{\vb{p}}
 \begin{pmatrix}-C_{\vb{p}}\eta_s\\\eta_s\end{pmatrix},
 \qquad H_0(\vb{p})W_s(\vb{p})=-E_{\vb{p}}W_s(\vb{p}),
 \label{eq:feqo-negative-hamiltonian-spinor}
\end{equation}
其中 $\eta_s$ 是二分量单位正交基。$\{U_1,U_2,W_1,W_2\}$ 是普通内积下的完整正交基。
\end{proposition}

\begin{proof}
Dirac 本征方程的上下两块为
\begin{equation*}
 (E-m_{\mathrm e}c^2)\phi=c\boldsymbol\sigma_{\rm P}\cdot\vb{p}\,\chi,
 \qquad
 (E+m_{\mathrm e}c^2)\chi=c\boldsymbol\sigma_{\rm P}\cdot\vb{p}\,\phi.
\end{equation*}
正能根给 $\chi=C_{\vb{p}}\phi$；负能根给 $\phi=-C_{\vb{p}}\chi$。因为 $C_{\vb{p}}$ 厄米且
\begin{equation*}
 C_{\vb{p}}^2=
 \frac{E_{\vb{p}}-m_{\mathrm e}c^2}{E_{\vb{p}}+m_{\mathrm e}c^2}I_2,
 \qquad N_{\vb{p}}^2(I_2+C_{\vb{p}}^2)=I_2,
\end{equation*}
有 $U_s^\dagger U_{s'}=W_s^\dagger W_{s'}=\delta_{ss'}$。不同能量支的正交性也直接给出
\begin{equation*}
 U_s^\dagger W_{s'}=N_{\vb{p}}^2\xi_s^\dagger(-C_{\vb{p}}+C_{\vb{p}})\eta_{s'}=0.
\end{equation*}
令 $h=H_0/E_{\vb{p}}$，则 $h=h^\dagger$、$h^2=I_4$，从而 $(I_4\pm h)^2=2(I_4\pm h)$ 和 $(I_4+h)(I_4-h)=0$ 给出所有投影恒等式。$\Tr H_0=0$ 又给 $\Tr\Pi_\pm=2$。因此
\begin{equation*}
 \Pi_+=\sum_sU_sU_s^\dagger,
 \qquad \Pi_-=\sum_sW_sW_s^\dagger.
\end{equation*}
对于\eqnrefs{eq:explicit-v}，$H_0(-\vb{p})v_s(p)=-E_{\vb{p}}v_s(p)$，故 $W_s(\vb{p})=\sqrt{c/(2E_{\vb{p}})}\,v_s(p^0,-\vb{p})$。负频反粒子旋量与同一三动量的负能哈密顿量旋量不能使用同一个动量实参。
\end{proof}

\begin{proposition}[协变自旋和与概率投影的换算]
未来指向在壳动量 $p$ 对应
\begin{align}
 \sum_sU_s(\vb{p})\bar U_s(\vb{p})
 &=\Pi_+(\vb{p})\gamma^0
 =\frac{c(\slashed p+m_{\mathrm e}c)}{2E_{\vb{p}}},
 \label{eq:feqo-probability-spin-sum}\\
 \sum_sW_s(\vb{p})\bar W_s(\vb{p})
 &=\Pi_-(\vb{p})\gamma^0
 =\frac{c(\slashed{\tilde p}-m_{\mathrm e}c)}{2E_{\vb{p}}},
 \qquad \tilde p=(E_{\vb{p}}/c,-\vb{p}).
 \label{eq:feqo-negative-probability-spin-sum}
\end{align}
普通内积的正交投影为 $\Pi_\pm$；含 Dirac 共轭的自旋和并非普通正交投影。
\end{proposition}

\begin{proof}
右乘 $\gamma^0$ 时，$\alpha_i\gamma^0=-\gamma^i$，所以
\begin{equation*}
 \Pi_+\gamma^0=\frac{E_{\vb{p}}\gamma^0-c\gamma^ip_i+m_{\mathrm e}c^2I_4}{2E_{\vb{p}}},
 \qquad
 \Pi_-\gamma^0=\frac{E_{\vb{p}}\gamma^0+c\gamma^ip_i-m_{\mathrm e}c^2I_4}{2E_{\vb{p}}}.
\end{equation*}
两式分别是 $c(\slashed p+m_{\mathrm e}c)/(2E_{\vb{p}})$ 和 $c(\slashed{\tilde p}-m_{\mathrm e}c)/(2E_{\vb{p}})$。同时
\begin{equation*}
 \bar U_sU_{s'}=\frac{m_{\mathrm e}c^2}{E_{\vb{p}}}\delta_{ss'},
 \qquad
 \bar W_sW_{s'}=-\frac{m_{\mathrm e}c^2}{E_{\vb{p}}}\delta_{ss'}.
\end{equation*}
协变归一化只需乘回 $\sqrt{2E_{\vb{p}}/c}$，得到\eqnrefs{eq:u-complete,eq:v-complete}。
\end{proof}

\begin{exbox}[$200\,\mathrm{keV}$ 电子的上下分量与投影]
取动能为 $200\,\mathrm{keV}$ 的电子，$\vb{p}=p\vb{e}_z$、$m_{\mathrm e}c^2=511\,\mathrm{keV}$。求该电子正能旋量的上下分量比、基矢的正能投影概率与流速。
\textbf{解。} 总能量是 $E=711\,\mathrm{keV}$，
\begin{equation*}
 pc=\sqrt{711^2-511^2}\,\mathrm{keV}=494.4\,\mathrm{keV},
 \qquad r=\frac{pc}{E+m_{\mathrm e}c^2}=0.4046,
 \qquad N^2=\frac{1222}{1422}=0.85935.
\end{equation*}
自旋向上时 $U_\uparrow=N(1,0,r,0)^{\mathsf T}$，因此上下两块的概率分别为 $N^2=0.85935$ 和 $N^2r^2=0.14065$；下分量已占约 $14\%$。对坐标基向量 $e_1=(1,0,0,0)^{\mathsf T}$，
\begin{equation*}
 \Pi_+e_1=N^2(1,0,r,0)^{\mathsf T},
 \qquad \|\Pi_+e_1\|^2=N^2.
\end{equation*}
所以 $e_1$ 并非纯正能电子；归一化其正能部分才得到 $U_\uparrow$。同一旋量的流速为 $c^2p/E=0.6953c$，与精确色散的群速度相同。
\end{exbox}

把正能旋量叠加成波包以后，概率流由同一个 Dirac 方程确定。
\label{sec:feqo-current-and-packet}
\begin{definition}[正能单电子波包与概率流]
以普通三维 $\delta$ 归一化动量基写入射电子，
\begin{equation}
 \Psi(\vb{r},t)=\sum_s\int\frac{\dd^3p}{(2\pi\hbar)^{3/2}}
 a_s(\vb{p})U_s(\vb{p})
 \ee^{\ii(\vb{p}\cdot\vb{r}-E_{\vb{p}}t)/\hbar},
 \qquad \sum_s\int\dd^3p\,|a_s(\vb{p})|^2=1.
 \label{eq:feqo-positive-wavepacket}
\end{equation}
概率密度、概率流和电荷四流依次为
\begin{equation}
 \rho=\Psi^\dagger\Psi,\qquad
 \vb{j}_{\rm prob}=c\Psi^\dagger\boldsymbol\alpha\Psi,\qquad
 j_{\rm ch}^\mu=q_{\mathrm e}c\bar\Psi\gamma^\mu\Psi
 =(q_{\mathrm e}c\rho,q_{\mathrm e}\vb{j}_{\rm prob}).
 \label{eq:feqo-dirac-current}
\end{equation}
$[\rho]=\mathrm{m^{-3}}$、$[\vb{j}_{\rm prob}]=\mathrm{m^{-2}s^{-1}}$；进入 Maxwell 方程的是电荷流。
\end{definition}

\begin{proof}[证明（归一化、连续性与群速度）]
对位置内积先积分平面波，得到 $(2\pi\hbar)^3\delta^{(3)}(\vb{p}-\vb{q})$。在分布支撑上 $U_s^\dagger(\vb{p})U_{s'}(\vb{p})=\delta_{ss'}$，因此
\begin{equation*}
 \int\dd^3r\,\Psi^\dagger\Psi
 =\sum_s\int\dd^3p\,|a_s(\vb{p})|^2=1.
\end{equation*}
将 Dirac 哈密顿量方程与其厄米共轭相加，质量项和实标势相消，微分动能给
\begin{equation*}
 \partial_t(\Psi^\dagger\Psi)
 =-c\big[(\partial_i\Psi^\dagger)\alpha_i\Psi
 +\Psi^\dagger\alpha_i\partial_i\Psi\big]
 =-\nabla\cdot\vb{j}_{\rm prob}.
\end{equation*}
波包衰减使无穷远通量消失，所以总概率守恒。对固定动量的正能旋量，Pauli 反对易关系直接给
\begin{align*}
 U_s^\dagger\alpha_iU_s
 &=N_{\vb{p}}^2\xi_s^\dagger(\sigma_iC_{\vb{p}}+C_{\vb{p}}\sigma_i)\xi_s\\
 &=N_{\vb{p}}^2\frac{2cp_i}{E_{\vb{p}}+m_{\mathrm e}c^2}
 =\frac{cp_i}{E_{\vb{p}}}.
\end{align*}
故 $\vb{j}_{\rm prob}/\rho=c^2\vb{p}/E_{\vb{p}}=\nabla_{\vb{p}}E$。波包中不同动量还存在干涉流，不能把该单平面波比值逐点代替完整电流。
\end{proof}
